\CWHeader{Лабораторная работа \textnumero 3}

\CWProblem{
Для реализации словаря из предыдущей лабораторной работы необходимо провести исследование скорости выполнения и потребления оперативной памяти. В случае выявления ошибок или явных недочётов, требуется их исправить.

Результатом лабораторной работы является отчёт, состоящий из:
\begin{itemize}
    \item Дневника выполнения работы, в котором отражено что и когда делалось, какие средства использовались и какие результаты были достигнуты на каждом шаге выполнения лабораторной работы.
    \item Выводов о найденных недочётах.
    \item Сравнения работы исправленной программы с предыдущей версией.
    \item Общих выводов о выполнении лабораторной работы, полученном опыте.
\end{itemize}

Минимальный набор используемых средств должен содержать утилиту \texttt{gprof} и библиотеку \texttt{dmalloc} (допускается замена на аналогичные или более развитые утилиты, такие как \texttt{Valgrind}).
}

\pagebreak